\section{Land, timberland and farmland}

\begin{qfmethod}[Raw land, timberland and farmland: their features, the sources of their returns, and analysing their risk, return and diversification]\label{meth:land-timber-farmland}
\qffield{Statement} \emph{Raw land} produces no income and costs property taxes and upkeep, so its return is appreciation alone. A holder needs price appreciation above the cost of capital plus the tax rate for a positive net present value. Alternatively the land is an option to develop, worth holding even when expected appreciation is lower, if land prices are volatile \citep[p.~753]{investment-valuation-2012}.

\emph{Timberland} returns come from biological growth of the trees, the price of timber, harvest income and the value of the land. Owners can delay harvests when prices are low, ``storing'' value on the stump. Its yield, like farmland's, depends on commodity prices rather than rents \citep[p.~143]{demaria-private-markets-2021}.

\emph{Farmland} earns crop income or cash rent and land appreciation. Values reflect returns per acre, interest rates and alternative investments, and parcel attributes such as soil, government payments, amenities and urban proximity. In 2026 US cropland averaged \$6,020 an acre with average cash rent of \$160; farm real estate was \$3.60 trillion, 83 percent of US farm assets \citep[ERS Farmland Value page, 2026 update]{nr-usda-farmland-2026}.

\emph{Risk, return and diversification.}
\begin{itemize}
\item Return sources: income yield, growth (biological or rent) and land appreciation; risks are commodity prices, weather and disease, fire, regulation and interest rates (Interpretation, from the sources above).
\item Both are offered as inflation hedges and diversifiers (Demaria, p.~143). In 2022 the NCREIF Timberland Index returned 12.90 percent, its best year since 2007, while property was weakening; in the first quarter of 2023 it returned 1.75 percent against $-1.81$ percent for the NCREIF Property Index \citep[p.~28]{ncreif-timberland-2023q1}.
\item Their published returns are appraisal-based, like NCREIF property and farmland series, and inherit smoothing \citep[pp.~742--743]{investment-valuation-2012}; \citep[pp.~416--421]{s7-ang}.
\item Access is through specialist funds, separate accounts or listed timber REITs; the private market is small, about 0.7 percent of private real assets for US timberland funds (Demaria, p.~143).
\end{itemize}

\qffield{Practice} \emph{Raw land} (Damodaran's example re-worked). At a 10 percent cost of capital and a 2 percent annual property tax on value, a one-year holding needs $(1+g)(1-0.02)=1.10$, or appreciation of 12.24 percent; Damodaran rounds this to 12 percent.

\emph{Farmland} (USDA averages, original costs). A gross rent yield of $160/6{,}020=2.66$ percent falls to 2.16 percent after costs of 0.5 percent of value. With rents growing 3 percent a year, the Gordon return is $2.16\%\times1.03+3\%=5.22$ percent: most of the return must come from growth.

\emph{Timber stand} (original). Four percent biological growth and a 1 percent rise in timber prices raise timber value by 5.04 percent without any market move. With a 3 percent harvest yield and 2 percent land appreciation the total is about 10 percent. Delaying the harvest in a weak market defers the yield but keeps the growth.

\emph{Diversification decision}: for an investor whose property holdings are equity-like and rate-sensitive, a small timberland allocation adds a return source (biological growth) unrelated to rents. The appraisal-based correlation should be discounted, and the low current farmland yield means the case rests on inflation and land growth.

\qffield{Conditions} USDA figures are national averages for 2026 from the ERS summary of NASS data; the timber and cost assumptions are illustrative.

\qffield{Pitfalls} Treating appraisal-based timber and farmland returns as risk measures. Ignoring the carrying cost of raw land. Assuming farmland rents rise with inflation every year.

\qffield{See also} \hyperref[meth:real-assets-allocation]{real asset allocation}; \hyperref[def:commodity-features]{commodities}.

\end{qfmethod}
